%% kapa: exam=1 round=25 year=2014 period=2 form=A subject=appraisal_law subject_ko=부동산관계법규 seq=40 number=40
%% source: Q-Net 감정평가사 기출문제 (cst00309) artlSeq=5101207 file=extracted/1008/5101207/2014년 제25회 감정평가사 1차시험 기출문제/2교시(A형).hwp
%% conversion: kordoc 4.24.51 -> tools/gen1.py (auto); review=proofread-sonnet; flags=none
\begin{question}{40}
부동산등기법령상 등기의 신청 등에 관한 설명으로 옳지 않은 것은?
\begin{choices}
  \item 등기원인이 발생한 후에 등기권리자에 대하여 상속이 있는 경우에는 상속인이 그 등기를 신청할 수 있다.
  \item 관리인이 있는 법인 아닌 재단에 속하는 부동산의 등기에 관하여는 그 재단을 등기권리자 또는 등기의무자로 한다.
  \item 채권자는 채무자를 대위하여 등기를 신청할 수 있다.
  \item 등기신청의 취하는 등기관이 등기를 마치기 전까지 할 수 있다.
  \item 같은 채권의 담보를 위하여 소유자가 다른 여러 개의 부동산(같은 등기소의 관할 내에 소재함)에 대한 저당권설정등기를 신청하는 경우에는 1건당 1개의 부동산에 관한 신청정보를 제공하는 방법으로 등기를 신청하여야 한다.
\end{choices}
\end{question}
\answer{5}
